
\begin{table}[htbp]
\centering
\caption{SOE Moderation Analysis}
\label{tab:soe_moderation}
\begin{threeparttable}
\begin{tabular}{lccc}
\toprule
 & \multicolumn{3}{c}{Dependent Variable: Stock Returns} \\
\cmidrule(lr){2-4}
 Variable & Full Sample & Non-SOEs & SOEs \\
 & (1) & (2) & (3) \\
\midrule
Constant & 4.000e-04*** & (1.000e-04) & 4.000e-04*** & (1.000e-04) & 5.000e-04*** & (1.000e-04) \\
Ambiguity Beta & -4.000e-04*** & (1.000e-04) & -4.000e-04*** & (1.000e-04) & -1.000e-04 & (1.000e-04) \\
Ambiguity $	imes$ SOE & 3.000e-04** & (1.000e-04) &  &  \\
Market Beta & 1.000e-04 & (1.000e-04) & 1.000e-04 & (1.000e-04) & 2.000e-04* & (1.000e-04) \\
SOE Dummy & 1.000e-04* & (0.000) &  &  \\
Log Market Cap & -1.000e-04* & (0.000) & -1.000e-04* & (0.000) & -1.000e-04 & (0.000) \\
Book-to-Market & 2.000e-04** & (1.000e-04) & 2.000e-04** & (1.000e-04) & 1.000e-04 & (1.000e-04) \\

\midrule
R-squared & 0.298 & 0.287 & 0.156 \\\\
Observations & 2,567 & 1,667 & 900 \\
\bottomrule
\end{tabular}
\begin{tablenotes}
\small
\item This table presents the results testing Hypothesis H4 that SOE status moderates the pricing of ambiguity. Column (1) includes an interaction term between ambiguity beta and SOE status. Columns (2) and (3) present subsample analysis for non-SOEs and SOEs, respectively. *, **, *** denote statistical significance at the 10\%, 5\%, and 1\% levels, respectively.
\end{tablenotes}
\end{threeparttable}
\end{table}
